\chapter{High-level overview}
For both the client and server, there exists only one entry point, which is either a local script or a server script. This means that the entry point must require other modules for core game logic. The requires are hard-coded, to be able to define order.
%--------------------------
\chapter{Shared Modules Overview}
Shared modules are both accessible to the server and client. They should be used for shared game logic, or shared data values. They should reside in the replicated storage.
%--------------------------
\chapter{Client System Overview}
Client modules should reside in StarterPlayerScripts. The client scripts are those, which the server has no use for.
Below is a high-level overview of the client sided architecture.
\section{High-Level Architecture}
\begin{itemize}
    \item \textbf{Input handler}
    \item \textbf{Character movement handler}
    \item \textbf{Client replication handler}        
\end{itemize}
\section{Dependencies}
\section{High-level diagram}

\begin{center}
\resizebox{\textwidth}{!}{
    \begin{tikzpicture}[]
    \node[startstop] (EntryPoint) {Entry Point};
    \node[process, right=of EntryPoint] (LoadPlayer) {Load player};
    \node[process, right=of LoadPlayer] (InputHandler) {Input handler};
    \node[process, right=of InputHandler] (Cmove) 
    {Character movement handler};
    \node[process, right=of Cmove] (Crep){Client replication handler};
    \node[process, right of=Crep]  (shandler){State handler};
    % Arrows
\end{tikzpicture}
}
\end{center}

%--------------------------
\chapter{Server System Overview}
Server modules should reside in ServerScriptService. Server modules are those modules, which would be a security risk for the client to have access to, or those modules which the client has no use for.
\section{High-Level Architecture}
\section{Dependencies}

%--------------------------
